\chapter{Panoramica dell'architettura}

\section{Struttura delle cartelle}

Il backend Ruggine è organizzato secondo una struttura modulare e chiara. Nella root del progetto si trova il file \texttt{Cargo.toml} e la cartella \texttt{src}, che contiene tutte le componenti principali del backend.



All'interno di \texttt{src} troviamo:

	extbf{\texttt{controllers}}: gestisce le richieste HTTP e WebSocket, occupandosi della logica di routing e dell'interazione con i servizi applicativi.

	extbf{\texttt{services}}: contiene la logica di business, implementando le regole applicative e coordinando le operazioni tra controller e repository.

	extbf{\texttt{repositories}}: si occupa dell'accesso ai dati e delle query verso il database, fornendo un'interfaccia astratta per la persistenza.

	extbf{\texttt{DTOs}}: definisce le strutture dati utilizzate per la comunicazione tra i vari livelli dell'applicazione e verso l'esterno.

	extbf{\texttt{migrations}}: contiene gli script SQL per la creazione e modifica delle tabelle del database.

I file principali di avvio e configurazione sono \texttt{error.rs}, \texttt{main.rs}, \texttt{router.rs} e \texttt{state.rs}.

\section{Componenti principali}

Il backend è sviluppato in Rust utilizzando le seguenti tecnologie:
Axum per la gestione del web server e delle rotte, SQLx per l'accesso asincrono al database, Tokio per la concorrenza e la gestione degli eventi, Tracing per il logging strutturato. Il frontend, invece, è realizzato in React e si trova nella cartella \texttt{ruggine\_client}.

\section{Flusso dati generale}

Il flusso dati tipico inizia con una richiesta HTTP o WebSocket ricevuta da un controller. Il controller valida e pre-elabora la richiesta, quindi delega la logica applicativa al service appropriato. Il service, a sua volta, interagisce con i repository per accedere o modificare i dati nel database. I DTOs vengono utilizzati per trasferire i dati tra i vari livelli dell'applicazione, garantendo coerenza e separazione delle responsabilità.